\documentclass[11pt]{article}

\usepackage[margin=1in]{geometry}
\usepackage[brazilian]{babel}
\usepackage{amsmath,amssymb,amsthm,mathtools}
\usepackage{enumitem}
\usepackage{hyperref}
\hypersetup{colorlinks=true,linkcolor=blue,citecolor=blue,urlcolor=blue}
\usepackage{orcidlink}
\usepackage{fontawesome5}

\newtheorem{claim}{Afirmação}[section]

\newcommand{\ZZ}{\mathbb{Z}}
\newcommand{\NN}{\mathbb{N}}

\title{Erros em Alegações de Prova (ou Refutação) Completa\\
da Conjectura de Collatz: Catálogo e Taxonomia\\
de Doze Casos}

\author{Renato Augusto Tavares~{\small\href{mailto:dr.renatotavares@gmail.com}{\textcolor{black}{\faEnvelope}}}~\orcidlink{0009-0002-0196-3311}\\
\normalsize Universidade Federal de Goiás}
\date{\today}

\begin{document}

\maketitle

\begin{abstract}
A Conjectura de Collatz atrai um volume incomum de alegações de prova e
refutação completas, majoritariamente veiculadas em servidores de
preprint e canais de baixo escrutínio, não em periódicos de matemática
com revisão por pares. Relatamos aqui uma revisão sistemática de doze
dessas alegações, encontradas ao longo de um levantamento bibliográfico
mais amplo, cada uma lida na íntegra, reconstruída de forma
independente e verificada computacionalmente sempre que a alegação
admite um teste numérico. Nenhuma das doze é uma prova ou refutação
válida. A contribuição desta nota não é a contagem em si, mas uma
taxonomia de como elas falham: duas categorias respondem por dez dos
doze casos (raciocínio circular que embute a conclusão da conjectura
numa definição ou premissa, três casos; generalização ilegítima de uma
única sequência adversária escolhida à mão para todas as sequências,
dois casos), e uma única falácia recorrente --- tratar uma propriedade
agregada, média ou assintótica como uma garantia determinística e
pontual para uma trajetória específica --- responde sozinha por cinco
casos, disfarçada respectivamente pelo Teorema de Baker, por uma
heurística de passeio aleatório, por um argumento de Lyapunov 2-ádico,
pela notação Big-O e pela densidade 2-ádica estática. Os dois casos
restantes são um erro elementar sem relação com a dificuldade real da
conjectura, e uma lacuna de rigor genuína que, ao contrário de seu
análogo mais próximo no catálogo, não produziu contraexemplo
computacional. Dez dos doze fracassos convergem, por rotas
superficialmente não relacionadas, para o mesmo obstáculo estrutural
que nenhum método conhecido resolve: a impossibilidade de limitar
sequências adversárias de passos crescentes sem, em algum ponto, supor
exatamente o que se quer provar.
\end{abstract}

\section{Introdução}
\label{sec:intro}

A Conjectura de Collatz está aberta há quase noventa anos, e seu
enunciado é simples o bastante para atrair um número incomum de
alegações de solução completa vindas de fora da comunidade de teóricos
dos números profissionais. Quase todas circulam por repositórios de
preprint, sites pessoais ou periódicos sem reputação estabelecida em
matemática pura, não por canais com revisão por pares especializada.
Nenhuma foi aceita pela comunidade matemática. Isso não é, por si só,
evidência de que uma alegação específica esteja errada --- mas
estabelece o prior correto: uma alegação de prova (ou refutação)
completa de um problema tão resistente, vinda de um canal sem revisão
por pares especializada, deve ser lida com a hipótese de trabalho de
que contém um buraco real, provavelmente exatamente no ponto em que o
``obstáculo estrutural'' conhecido da conjectura precisaria ser
superado (a impossibilidade, com qualquer método conhecido hoje, de
limitar o quão longa pode ser uma sequência adversária de passos
crescentes --- ver \cite{LagariasSurvey} para o relato padrão de por
que isso é difícil).

Esta nota cataloga doze dessas alegações, encontradas ao longo de um
levantamento bibliográfico mais amplo conduzido paralelamente a um
programa de pesquisa não relacionado sobre uma generalização $qx+1$ da
medida de Syracuse de Tao \cite{Tao2022}. Cada alegação foi lida na
íntegra, seu argumento central reconstruído de forma independente, e
verificada computacionalmente sempre que admite um teste numérico;
vários dos julgamentos mais delicados também foram cruzados com um
revisor matemático independente baseado em modelo de linguagem. Nenhuma
das doze sobrevive ao escrutínio como prova ou refutação válida.

O ponto desta nota não é o fato simples de que doze alegações amadoras
estão erradas --- isso seria pouco notável e pouco interessante por si
só. O ponto é \emph{como} elas falham. A Seção~\ref{sec:taxonomy} agrupa
os doze casos em cinco categorias pela forma de seu erro central, com
um resumo breve de cada caso; três categorias, cobrindo dez casos,
convergem para o mesmo obstáculo estrutural por rotas que parecem, à
primeira vista, inteiramente não relacionadas entre si. A
Seção~\ref{sec:discussion} discute o que essa convergência sugere sobre
por que a conjectura resiste a esses estilos particulares de ataque, e
o que as duas exceções ao padrão revelam por contraste. A
Seção~\ref{sec:conclusion} conclui.

Para deixar claro o que esta nota \emph{não} é: não é uma alegação
sobre a competência ou a boa-fé dos autores, vários dos quais são
explícitos e cuidadosos sobre os limites de seus próprios argumentos em
outras partes do mesmo texto (rotulando corretamente um passo em aberto
como ``conjectura'', não ``teorema'', por exemplo). O foco em todo o
texto é estritamente o conteúdo matemático do erro específico
identificado em cada caso.

\section{Metodologia}
\label{sec:methodology}

Cada uma das doze alegações foi: (i) lida na íntegra, não apenas
resumo e conclusão; (ii) reconstruída como uma cadeia explícita de
premissas e inferências, de modo a isolar exatamente o passo que carrega
o peso lógico de ``prova'' ou ``refuta''; (iii) verificada
computacionalmente sempre que a alegação, ou o passo suspeito de ser o
buraco, admite um teste numérico --- seja por reimplementação direta da
identidade alegada, seja pela construção de um contraexemplo dirigido;
e (iv), nos poucos casos em que o julgamento matemático envolvido era
sutil o bastante para justificar uma segunda opinião, cruzada com um
revisor independente baseado em modelo de linguagem, sem acesso prévio
à nossa própria conclusão de trabalho. O detalhe técnico completo de
cada caso --- o argumento reconstruído exato e a verificação
computacional --- foi registrado durante esta revisão em notas de
trabalho internas; esta nota resume esse material, não o reproduz por
completo.

\section{Taxonomia dos padrões de erro recorrentes}
\label{sec:taxonomy}

Uma primeira tentativa desta taxonomia, montada apenas a partir do
título e de uma linha de status de cada caso, propôs sete categorias.
Ela não sobreviveu ao contato com os textos completos: um caso (Boyle,
\S\ref{sec:catC}) não se encaixava em nenhuma das sete, e a leitura
integral dos casos restantes mostrou que ele, junto com mais quatro,
era na verdade o mesmo erro vestido em aparato formal não relacionado.
A taxonomia abaixo reflete a leitura completa e tem cinco categorias,
uma das quais --- Seção~\ref{sec:catC} --- responde sozinha por cinco
dos doze casos.

\subsection{Categoria A: raciocínio circular}
\label{sec:catA}

O objeto central do argumento, ou seu passo central, só está bem
definido ou é válido sob a suposição de que a conjectura já vale --- de
modo que a ``prova'' não estabelece a terminação, apenas a reafirma em
outra notação. Três casos:

\textbf{Getachew (2025)} \cite{Getachew2025} constrói uma árvore reversa
do mapa de Collatz e define uma relação ``pai'' que se revela, termo a
termo, idêntica ao próprio mapa direto de Collatz
($\mathrm{pai}(m) = m/2$ para $m$ par, $3m+1$ para $m$ ímpar). O
Lema~4.3 do paper, ``todo caminho de volta à raiz é finito'', é então
logicamente equivalente à própria conjectura, não uma consequência da
aciclicidade da árvore e da unicidade do pai, como se alega:
aciclicidade e pai único garantem que um caminho nunca repete nem
bifurca, nunca que seja finito. O teorema de cobertura do paper
(Teorema~5.1) falha por um motivo paralelo: prova uma decomposição
aritmética universal (número ímpar vezes potência de 2) que
confirmamos computacionalmente nunca usar de fato a regra de
ramificação.

\textbf{Spencer (2025a)} \cite{Spencer2025a} constrói uma árvore reversa
diferente, indexada por classes residuais módulo $2\cdot 3^J$, com um
argumento de contador ternário parcialmente formalizado em Lean. Prova
corretamente que toda classe residual primitiva está ``ocupada'' em
toda escala finita $J$ --- mas seu Teorema~15.3 conclui então que todo
inteiro específico tem um endereço reverso finito, um passo nunca
justificado: ``resíduo ocupado'' não implica ``valor exato atingido''.
Confirmamos essa lacuna diretamente: a classe residual de $n=27$ está
ocupada na árvore real construída até profundidade 6 e 8, mas apenas
por representantes da ordem de $10^6$--$10^8$, nunca pelo próprio 27.

\textbf{Yun (2026)} \cite{Yun2026} contém três argumentos circulares
independentes, cada um dos quais, sozinho, já invalida a alegação do
paper. Sua ``função de posto'' central
$r(x) := \min\{n \in \NN : f^{(n)}(x) = 1\}$ só está definida para
valores de $x$ cuja órbita já atinge 1 --- se alguma órbita nunca
atingisse 1, exatamente a possibilidade que a conjectura precisa
excluir, $r(x)$ simplesmente não existiria, de modo que o argumento de
descida bem-fundada construído sobre ela já pressupõe sua conclusão.
Confirmamos isso aplicando a mesma construção a $g(x)=2x$, um mapa que
provadamente diverge para todo $x>0$: a função de posto análoga fica
indefinida para os 18 valores testados, exatamente como aconteceria com
$f$ se alguma órbita não atingisse 1. Separadamente, o argumento de
exclusão de ciclos do paper (Teorema~6.1) afirma que ``o único ponto
fixo conhecido é $x=1$'' como se isso excluísse outros ciclos, e seu
argumento de ``compressão informacional'' (\S5, \S9.2.3) é um non
sequitur sobre cardinalidades de conjuntos infinitos, compatível com
qualquer comportamento de longo prazo, não especificamente com a
convergência a 1.

\subsection{Categoria B: generalização ilegítima de uma única sequência}
\label{sec:catB}

O argumento mostra, corretamente, que uma sequência específica e
escolhida à mão de passos crescentes e decrescentes tem certa
propriedade, e então estende a conclusão para todas as sequências sem
limitar o quão longa uma sequência adversária de passos crescentes
poderia ser --- exatamente o passo que o obstáculo estrutural torna
difícil. Dois casos:

\textbf{Santos (2018)} \cite{Santos2018} introduz uma notação binária de
``movimentos''/``reduções fortes'' equivalente à sequência ordinária de
valuações 2-ádicas de uma órbita de Collatz. Seu argumento central
(\S2.6) escolhe um $K>30$ hipotético, \emph{assume uma sequência fixa e
específica} de passos (dois movimentos crescentes seguidos de três
decrescentes, com divisores escolhidos à mão), mostra algebricamente
que $K$ decresce ao longo \emph{daquela sequência particular}, e então
generaliza isso para todo $K>30$ sem nunca mostrar que sequências
adversárias de passos crescentes não podem ser arbitrariamente longas.

\textbf{Halemane (2025), o teorema CTUHSK} \cite{Halemane2025}
reformula a árvore reversa padrão com notação nova
(Binary-Exponential-Ladders, operadores $D^\#/U^\#$), mas não acrescenta
nada estruturalmente novo. Sua condição necessária é uma tautologia (o
conjunto \emph{definido} como quem alcança o ciclo trivial de fato
alcança o ciclo trivial); sua condição suficiente, que carrega o peso
real de excluir ciclos extras e cadeias divergentes, tem um furo
decisivo: o argumento de redução ao absurdo (\S10.2.1) assume que o
predecessor de um elemento mínimo de um ciclo hipotético é dado pela
fórmula do paper com expoente $v=1$ especificamente, quando essa
fórmula tem de fato infinitas soluções válidas, uma por expoente ímpar
$v=1,3,5,\dots$; apenas $v=1$ dá um valor menor (violando a
minimalidade suposta), enquanto todo $v\ge 3$ dá um valor maior, sem
gerar contradição alguma. Confirmamos computacionalmente que isso é um
padrão algébrico geral, não um acidente do exemplo numérico escolhido
pelo paper.

\subsection{Categoria C: propriedade agregada ou assintótica confundida com garantia determinística e pontual}
\label{sec:catC}

Esta categoria é o achado empírico central do paper. Cinco argumentos
superficialmente não relacionados --- construídos, respectivamente,
sobre o Teorema de Baker, a heurística clássica de passeio aleatório,
uma cota de Lyapunov 2-ádica, a notação assintótica Big-O e a
densidade 2-ádica estática --- cometem, na verdade, o erro idêntico:
uma propriedade que só vale como média, expectativa de ensemble ou cota
assintótica é substituída, sem justificativa, por uma garantia que
precisaria valer deterministicamente, em todo ponto, ao longo de uma
trajetória específica. Três dos cinco documentos-fonte abaixo já se
citam mutuamente quanto a esse padrão exato, independentemente deste
catálogo, o que sugere que se trata de um atrator natural para
tentativas de prova de Collatz, não um artefato da nossa classificação.

\textbf{Mohammed} \cite{Mohammed2024} constrói um sieve de densidade
geométrica correto e exclui corretamente 1-ciclos, então tenta excluir
ciclos não-triviais via o Teorema de Baker sobre formas lineares em
logaritmos. A equação de ciclo de fato força $M/P \to \log_2 3$ para um
ciclo real e autoconsistente com elementos grandes --- mas o paper
substitui $M \approx 2P$, a expectativa \emph{ergódica, média} para um
inteiro ímpar genérico (já conhecida neste projeto como fato de
referência, não como restrição de ciclo), produzindo uma checagem que
não restringe nada.

\textbf{Boyle (2026)} \cite{Boyle2026} deriva, via a heurística padrão
de passeio aleatório, que a fração média de ensemble de passos pares ao
longo de órbitas de Collatz é $2/3$ --- válida, na melhor das
hipóteses, como média sobre muitos valores iniciais $n$ aleatórios ---
e então substitui essa constante diretamente numa equação (Lema~4.2)
sobre uma trajetória hipotética específica e fixa, derivando uma
``contradição'' que depende inteiramente dessa substituição ilegítima.
Confirmamos computacionalmente que, mesmo entre órbitas conhecidas por
convergir, a fração de passos pares realizada por cada trajetória
individual varia substancialmente ao redor de $2/3$ (média 0,677,
desvio-padrão 0,031 sobre 3000 trajetórias amostradas) --- é uma
estatística agregada, não uma restrição sobre qualquer órbita
individual.

\textbf{Tynski (2026)} \cite{Tynski2026} constrói um aparato formal
extenso e majoritariamente correto (álgebra de Clifford, um traço
explícito, identidades de exclusão de ciclo já estabelecidas) em torno
de um ``teorema mestre'' que unifica a Hipótese de Riemann e a
Conjectura de Collatz, mas a exclusão de órbitas divergentes repousa
inteiramente sobre um axioma que afirma uma cota de Lyapunov
\emph{determinística}
$\log_2 n_m \le \log_2 n_0 - m\delta + C\sqrt m$ para uma constante
universal $C$, justificada de forma circular (o texto deriva o
decréscimo necessário ``uma vez que $A_m/m > \log_2 3$ seja forçado'',
exatamente a propriedade que estaria sendo argumentada) e sustentada
apenas por verificação numérica até $n_0 \le 3\times 10^4$. Testamos a
alegação diretamente além dessa faixa: o $C$ mínimo necessário já
excede a cota universal alegada de 5 para sementes
$10^4 \le n_0 < 10^9$ (pior caso 5,08), e recordistas conhecidos de
atraso longo exigem $C$ até 8,34 --- uma refutação computacional
direta, não apenas ausência de prova.

\textbf{A disputa Syzdykov (2025--2026)}
\cite{Syzdykov2025,Syzdykov2026,LafontaineCheong2026} é uma troca de
três textos, não uma única alegação. A nota original \cite{Syzdykov2025}
afirma $O(f(n)) \le n$ para a função de Collatz $f$, substitui
$n = 2^m-1$, e declara $O(f(2^m-1)) \approx 3^m > 2^m - 1$ uma
contradição que refutaria a conjectura. Lafontaine \& Cheong
\cite{LafontaineCheong2026} identificam corretamente o erro categorial:
$O(\cdot)$ denota um conjunto de funções, não um escalar, de modo que
$O(f(n)) \le n$ não é bem-formado; e mesmo sob uma leitura pontual
caridosa, uma cota assintótica (``para algum $C,n_0$, para todo
$n \ge n_0$'') não licencia conclusão alguma sobre o valor exato de um
$n$ específico --- a mesma confusão entre uma cota existencial e
agregada e uma cota pontual e determinística vista nos outros quatro
casos desta categoria. A réplica de Syzdykov \cite{Syzdykov2026} não
rebate esse ponto; redefine $O(f(n))$ a posteriori para significar um
comprimento de órbita escalar e reafirma a conclusão original como
``já provada''.

\textbf{Barghout (2019)} \cite{Barghout2019} prova corretamente, por
indução elementar, que a valuação 2-ádica de inteiros pares (ou de
valores gerados por $3n+1$) segue um padrão de densidade fixo e
determinístico --- um fato estático sobre o \emph{conjunto} de inteiros
pares, não sobre nenhuma órbita específica. Seu argumento de
convergência então importa essa densidade estática como se fosse a
\emph{estatística de visitas} de uma órbita dinâmica específica,
assumindo implicitamente que a órbita equidistribui --- exatamente o
conteúdo não provado da conjectura, não uma consequência dos lemas de
densidade. Confirmamos por simulação direta que órbitas individuais
(por exemplo, $n=9663$) podem crescer a cerca de 2800 vezes seu valor
inicial antes de eventualmente decrescer, ilustrando concretamente por
que um argumento ``em média'' não restringe nenhuma trajetória
individual.

\subsection{Categoria D: um erro elementar sem relação com o obstáculo estrutural}
\label{sec:catD}

\textbf{Roif (2026)} \cite{Roif2026} constrói um espaço topológico de
três pontos $S=\{0,2,1\}$, reinterpreta o inteiro 2 como um
``recipiente vazio'' $\{\}$, e prova um lema afirmando que todo
subconjunto dos inteiros positivos com densidade assintótica zero deve
ser vazio. Isso é falso como fato matemático geral e é o único erro do
catálogo sem relação com a dificuldade real da conjectura: densidade
zero não implica finitude, muito menos vazio. O conjunto dos quadrados
perfeitos, ou das potências de 2, tem densidade zero e é claramente
infinito. Esse único erro elementar derruba todo o argumento de
fechamento do paper (seu Teorema~9.1 e Corolário~9.2, que de outra
forma completariam a prova), independentemente do restante do aparato
formal (identidades próximas de álgebra de Clifford, uma função de onda
tipo Lyapunov), que é uma reformulação correta, mas por si só
insuficiente, de resultados já padrão na literatura.

\subsection{Categoria E: uma lacuna de rigor genuína sem contraexemplo computacional}
\label{sec:catE}

\textbf{Spencer (2026)} \cite{Spencer2025b}, do mesmo autor do caso da
Categoria A na \S\ref{sec:catA}, usa uma árvore reversa diferente (grau
infinito, via o inverso do mapa acelerado). Diferente de seu paper
irmão, construir a árvore real e testar cobertura até $10\,000$ não
encontrou \emph{nenhuma} lacuna na prática --- aparentes falhas de
cobertura em orçamentos computacionais menores convergiam exatamente à
raiz quando checadas sem limite de magnitude. Ainda assim, identificamos
uma lacuna genuína no argumento escrito: o Teorema~14.1 prova ocupação
de \emph{classes} residuais, e o Teorema~14.2 conclui ocupação de
\emph{elementos} individuais sem justificar a passagem --- a mesma forma
lógica do caso da Categoria A, mas que, até agora, não se manifestou
como falha real sob busca computacional. Essa distinção importa: nem
toda lacuna numa tentativa de prova de Collatz é igualmente fatal, mesmo
dentro de argumentos intimamente relacionados do mesmo autor.

\section{Discussão}
\label{sec:discussion}

Dez dos doze casos --- todos das Categorias A, B e C --- falham, por
rotas que parecem não relacionadas à primeira vista, exatamente no
mesmo ponto: nenhum deles estabelece uma cota genuína sobre o quão
longa pode ser uma sequência adversária de passos crescentes, o
obstáculo estrutural que nenhum método conhecido resolve. A
Categoria~C sozinha responde por cinco dos dez, e é a parte mais
marcante do padrão: o mesmo erro conceitual --- substituir uma
propriedade agregada, média ou assintótica por uma garantia pontual e
determinística --- sobrevive vestido em cinco vocabulários formais não
relacionados (teoria dos números elementar via o Teorema de Baker, a
heurística clássica de passeio aleatório, uma cota de Lyapunov
2-ádica, a notação assintótica Big-O e a densidade 2-ádica estática).
Isso não é um artefato da nossa própria classificação: três dos cinco
documentos-fonte já se citam mutuamente, e dois deles citam
explicitamente a heurística de passeio aleatório, como instâncias da
mesma confusão, independentemente deste catálogo.

Os dois casos fora desse padrão são informativos por contraste. O erro
de Roif (Categoria D) é um engano isolado e elementar, sem relação com
a dificuldade real da conjectura --- um lembrete de que nem todo erro
numa tentativa de prova de Collatz é evidência sobre \emph{por que} a
conjectura é difícil, alguns simplesmente estão errados por motivos
não relacionados. O segundo paper de Spencer (Categoria E) mostra a
lição oposta: uma forma lógica idêntica a um fracasso da Categoria A
(ocupação de classe residual confundida com ocupação de elemento
individual) pode, num caso, não se manifestar como contraexemplo
computacional real --- reforçando que a taxonomia classifica a
\emph{estrutura} do argumento, não sua gravidade, e que lacunas
estruturalmente semelhantes nem sempre são igualmente fatais.

Notamos, por fim, que a coleção mais ampla de papers relacionados a
Collatz levantada junto com esses doze (fora do escopo desta nota)
contém muitos papers que não alegam prova completa alguma, e que são
cuidadosos em rotular passos em aberto como conjecturais. Os doze casos
aqui foram selecionados precisamente por fazerem a alegação forte; a
taxonomia acima não deve ser lida como sugerindo que papers de alegação
de prova são típicos da literatura mais ampla sobre a conjectura.

\section{Conclusão}
\label{sec:conclusion}

As doze alegações revisadas de prova ou refutação completa da
Conjectura de Collatz falham sob escrutínio independente, e nenhuma
está perto de ser válida. O valor deste catálogo não é a contagem, mas
a concentração: metade dos doze casos, apesar de usar vocabulários
matemáticos inteiramente diferentes, comete o erro idêntico de
confundir uma propriedade agregada ou assintótica com uma garantia
determinística e pontual, e mais três cometem um de dois erros
intimamente relacionados (circularidade, generalização ilegítima de
uma única sequência) que convergem para o mesmo obstáculo estrutural
conhecido. Essa regularidade é, em si, uma pequena evidência sobre por
que a conjectura resiste a ataques elementares: os modos de fracasso
não são diversos, eles se repetem.

\section*{Agradecimentos}

Partes da verificação descrita nesta nota (reconstrução de casos,
checagens computacionais e cruzamento de julgamentos matemáticos mais
sutis) foram realizadas com o auxílio de Claude (modelos das famílias
Claude 4 e 5, Anthropic), usado como assistente de pesquisa sob a
direção e revisão do autor.

\begin{thebibliography}{99}

\bibitem{LagariasSurvey} J.\ C.\ Lagarias (ed.), \emph{The Ultimate
  Challenge: The $3x+1$ Problem}, American Mathematical Society, 2010.

\bibitem{Tao2022} T.\ Tao, \emph{Almost all orbits of the Collatz map
  attain almost bounded values}, Forum of Mathematics, Pi \textbf{10}
  (2022), e12.

\bibitem{Getachew2025} E.\ S.\ Getachew, \emph{Unfolding the Collatz
  Tree: An Indirect Structural Proof}, Research in Mathematics
  (Taylor \& Francis) \textbf{12}(1), 2025.

\bibitem{Spencer2025a} M.\ E.\ Spencer, \emph{Finite Block Exhaustion
  and Rooted Occupancy in the Inverse Collatz System: A No-Alternative
  Proof via Ternary Counters, Primitive Carry, and Forced Thread
  Displacement}, academia.edu preprint, 2025.

\bibitem{Spencer2025b} M.\ E.\ Spencer, \emph{Rooted Surjectivity from
  the Invariant E/O Refinement System}, academia.edu preprint, 2026.

\bibitem{Santos2018} O.\ de O.\ Santos, \emph{Proving the Collatz
  Conjecture with Binaries Numbers}, Pure and Applied Mathematics
  Journal \textbf{7}(5) (2018), Science Publishing Group.

\bibitem{Halemane2025} K.\ P.\ Halemane, \emph{Collatz-Thwaites-Ulam-Hasse-Syracuse-Kakutani
  (CTUHSK) Theorem: Convergence of Collatz ($3n+1$) Sequence to the
  Trivial Cycle Proved}, engrXiv preprint, 2025.

\bibitem{Mohammed2024} A.\ Mohammed, \emph{Structural Analysis, Dynamic
  Density Sieve, and Logarithmic Contraction of Collatz Sequences},
  preprint, Faculty of Science, Kafr El-Sheikh University, 2026.

\bibitem{Boyle2026} D.\ G.\ Boyle, \emph{The Collatz Conjecture is
  True}, rxiverse.org preprint, 2026.

\bibitem{Roif2026} A.\ Roif, \emph{On the Convergence of the Collatz
  Function via Galois Fields, Topological Duality, and Lyapunov Wave
  Analysis}, academia.edu preprint, v4, 2026.

\bibitem{Yun2026} Y.\ H.\ Yun, \emph{A Structural Proof of the Collatz
  Conjecture via Non-Repeating Trajectory and Recursive Decay}, OSF
  preprint, Zero Theoretical Physics Lab, 2026.

\bibitem{Tynski2026} K.\ Tynski, \emph{A Common Proof of the Riemann
  Hypothesis and the Collatz Conjecture}, academia.edu preprint, 2026.

\bibitem{Syzdykov2025} M.\ Syzdykov, \emph{Disproof of Collatz
  Conjecture by Using O-notation}, ResearchGate preprint, 2025.

\bibitem{Syzdykov2026} M.\ Syzdykov, \emph{Continued Disproof Sentence
  towards Collatz Conjecture}, preprint, 2026.

\bibitem{LafontaineCheong2026} M.\ Lafontaine e G.\ Cheong, \emph{A
  Note on a Claimed Disproof of the Collatz Conjecture}, ResearchGate
  preprint, 2026.

\bibitem{Barghout2019} K.\ Barghout, \emph{On the Probabilistic Proof
  of the Convergence of the Collatz Conjecture}, Journal of Probability
  and Statistics (Hindawi) \textbf{2019}, Article ID 6814378,
  \url{https://doi.org/10.1155/2019/6814378}.

\end{thebibliography}

\end{document}
